\chapter{Introduction}
TODO

\section{Notation}
This section presents the notational conventions used throughout this document.

\subsection{Keywords}
Important terms and concepts that are introduced for the first time are
presented in an italic style, e.g. \emph{subject specification}. Subsequent
occurences of the same term have no special formatting.

\subsection{Numbers}
Regular numbers that have no leading special character are expressed as decimal
values. Hexadecimal numbers such as memory addresses are explicitly preceded by
\texttt{0x}.

Storage units such as kilo-, mega- and gigabyte are designated by the common
abbreviations KB, MB and GB.

\section{Related literature}
Since the target hardware platform of the separation kernel is the Intel x86
architecture, its specification called "Intel\textsuperscript{\textregistered}
64 and IA-32 Architectures Software Developer's Manual" \cite{IntelSDM} is the
main source of technical information about the hardware platform. The documents
are commonly referred by their short name Intel SDM.

The books are available online and updated by Intel on a regular basis. This
can lead to changes in the document structure. The chapter and section
citations given in this report refer to the Intel SDM revision 44, released in
August 2012.

\section{Provenance of project name}
\emph{Muen} is Japanese and translates to 'unrelated/without relation'. It was
chosen since it is a fitting allegory of the components executing on top of the
separation kernel. Figure \ref{fig:muen} depicts the word in Japanese kanji
characters.

\begin{figure}[h]
	\centering
	\includegraphics[scale=0.4]{images/muen.pdf}
	\caption{Muen in kanji}
	\label{fig:muen}
\end{figure}

The root of the word Muen is 'Mu' which denotes a negative: the absence of
something. It is a keyword in Chan and Zen Buddhism and also part of the Jargon
File \cite{jargonfile}.
